% =====================================================================
% comum/packages.tex — PONTO DE ENTRADA DOS MÓDULOS COMPARTILHADOS
% =====================================================================
% Uso (no preâmbulo do documento, logo após \documentclass):
%
%     % (opcional) configurações ANTES do \input:
%     \newcommand{\comumNumeracao}{chapter}   % chapter | section | (vazio)
%     \newcommand{\comumBibliografia}{referencias.bib}
%     % =====================================================================
% comum/packages.tex — PONTO DE ENTRADA DOS MÓDULOS COMPARTILHADOS
% =====================================================================
% Uso (no preâmbulo do documento, logo após \documentclass):
%
%     % (opcional) configurações ANTES do \input:
%     \newcommand{\comumNumeracao}{chapter}   % chapter | section | (vazio)
%     \newcommand{\comumBibliografia}{referencias.bib}
%     % =====================================================================
% comum/packages.tex — PONTO DE ENTRADA DOS MÓDULOS COMPARTILHADOS
% =====================================================================
% Uso (no preâmbulo do documento, logo após \documentclass):
%
%     % (opcional) configurações ANTES do \input:
%     \newcommand{\comumNumeracao}{chapter}   % chapter | section | (vazio)
%     \newcommand{\comumBibliografia}{referencias.bib}
%     % =====================================================================
% comum/packages.tex — PONTO DE ENTRADA DOS MÓDULOS COMPARTILHADOS
% =====================================================================
% Uso (no preâmbulo do documento, logo após \documentclass):
%
%     % (opcional) configurações ANTES do \input:
%     \newcommand{\comumNumeracao}{chapter}   % chapter | section | (vazio)
%     \newcommand{\comumBibliografia}{referencias.bib}
%     \input{comum/packages}
%
% Carrega, NESTA ORDEM, os módulos abaixo. A ordem importa:
%   1. idioma       — codificação, fontes e babel (sempre primeiro)
%   2. matematica   — ams*, mathtools, bm, tikz-cd
%   3. figuras      — graphicx, tikz, pgfplots, float, (sub)caption
%   4. layout-abnt  — margens ABNT, espaçamento 1,5, cabeçalho
%   5. referencias  — biblatex-abnt, hyperref e cleveref
%                     (hyperref/cleveref precisam vir depois de quase tudo)
%   6. theoremstyles — ambientes de teorema + nomes do cleveref
%                     (precisa do cleveref já carregado)
%   7. comandos     — macros próprias (\asyimage, ambiente citacao)
%
% Para desligar um módulo, comente a linha correspondente.
% Em beamer NÃO use este arquivo inteiro: beamer já define teoremas e
% layout; carregue só os módulos 1, 2, 3 e 5 (veja templates/slides-unespar).
% =====================================================================

\input{comum/pacotes/idioma}
\input{comum/pacotes/matematica}
\input{comum/pacotes/figuras}
\input{comum/pacotes/layout-abnt}
\input{comum/pacotes/referencias}
\input{comum/theoremstyles}
\input{comum/comandos}

%
% Carrega, NESTA ORDEM, os módulos abaixo. A ordem importa:
%   1. idioma       — codificação, fontes e babel (sempre primeiro)
%   2. matematica   — ams*, mathtools, bm, tikz-cd
%   3. figuras      — graphicx, tikz, pgfplots, float, (sub)caption
%   4. layout-abnt  — margens ABNT, espaçamento 1,5, cabeçalho
%   5. referencias  — biblatex-abnt, hyperref e cleveref
%                     (hyperref/cleveref precisam vir depois de quase tudo)
%   6. theoremstyles — ambientes de teorema + nomes do cleveref
%                     (precisa do cleveref já carregado)
%   7. comandos     — macros próprias (\asyimage, ambiente citacao)
%
% Para desligar um módulo, comente a linha correspondente.
% Em beamer NÃO use este arquivo inteiro: beamer já define teoremas e
% layout; carregue só os módulos 1, 2, 3 e 5 (veja templates/slides-unespar).
% =====================================================================

% =====================================================================
% comum/pacotes/idioma.tex — codificação, fontes e idioma
% =====================================================================
% Detecta o motor de compilação:
%   * XeLaTeX / LuaLaTeX (padrão do repositório, ver .vscode/settings.json):
%       fontspec com Latin Modern (Unicode nativo, sem inputenc).
%   * pdfLaTeX: inputenc utf8 + fontenc T1 + lmodern.
% Para trocar a fonte principal em XeLaTeX, use DEPOIS deste módulo:
%       \setmainfont{Liberation Serif}   % ex.: Times-compatível (IC/PRPPG)
%       \setsansfont{texgyreheros}[Extension=.otf, ...]  % Helvetica (slides)
% =====================================================================

\usepackage{iftex}

\ifPDFTeX
  \usepackage[utf8]{inputenc}
  \usepackage[T1]{fontenc}
  \usepackage{lmodern}
\else
  \usepackage{fontspec}
\fi

\usepackage[brazilian]{babel}

% Listas: enumitem (rótulos/espaçamento) e moreenum (\roman, \greek, ...)
\usepackage{enumitem}
\usepackage{moreenum}

% =====================================================================
% comum/pacotes/matematica.tex — pacotes de matemática
% =====================================================================
% amsmath/amssymb/amsthm : base AMS (amsthm é usado por theoremstyles.tex)
% mathtools              : extensões do amsmath (\coloneqq, \DeclarePairedDelimiter)
% bm                     : negrito em símbolos matemáticos (\bm{\alpha})
% amscd / tikz-cd        : diagramas comutativos (simples / completos)
% =====================================================================

\usepackage{amsmath}
\usepackage{amssymb}
\usepackage{amsthm}
\usepackage{mathtools}
\usepackage{bm}
\usepackage{amscd}
\usepackage{tikz-cd}

% =====================================================================
% comum/pacotes/figuras.tex — figuras, TikZ e legendas
% =====================================================================
% graphicx            : \includegraphics (imagens ficam em ../figures/)
% tikz + bibliotecas  : as usadas nos projetos (setas, ângulos, sombras...)
% pgfplots            : gráficos de funções
% float               : opção [H] para fixar figuras
% caption/subcaption  : legendas e subfiguras
% asymptote           : figuras 3D em .asy (TCC). Desligado por padrão
%                       porque exige o programa `asy` instalado; para
%                       ligar, defina \newcommand{\comumAsymptote}{} antes
%                       de carregar os módulos (veja comandos.tex: \asyimage).
% =====================================================================

\usepackage{graphicx}
\graphicspath{{../figures/}{figures/}}

\usepackage{tikz}
\usetikzlibrary{arrows.meta, calc, positioning, decorations.markings,
  angles, quotes, shapes.geometric, shapes.symbols, shadows}

\usepackage{pgfplots}
\pgfplotsset{compat=1.18}

\usepackage{float}
\usepackage{caption}
\usepackage{subcaption}

\ifdefined\comumAsymptote
  \usepackage{asymptote}
\fi

% =====================================================================
% comum/pacotes/layout-abnt.tex — página no padrão ABNT (NBR 14724)
% =====================================================================
% Margens: superior/esquerda 3cm, inferior/direita 2cm.
% Espaçamento 1,5; recuo no primeiro parágrafo; número da página no
% canto superior direito (fancyhdr); tabelas com booktabs/multirow.
%
% Documentos com cabeçalho/margens próprios (ex.: relatório PRPPG)
% podem sobrescrever depois com \geometry{...} e \fancyhead{...}.
% NÃO usar em beamer.
% =====================================================================

\usepackage[a4paper, top=3cm, bottom=2cm, left=3cm, right=2cm,
  headheight=14.5pt]{geometry}

\usepackage{setspace}
\onehalfspacing

\usepackage{indentfirst}
\setlength{\parindent}{1.25cm}

\usepackage{fancyhdr}
\pagestyle{fancy}
\fancyhf{}
\fancyhead[R]{\thepage}
\renewcommand{\headrulewidth}{0.4pt}

\usepackage{booktabs}
\usepackage{multirow}

% =====================================================================
% comum/pacotes/referencias.tex — bibliografia e referências cruzadas
% =====================================================================
% biblatex  : backend biber, estilo ABNT (pacote biblatex-abnt).
%             Arquivo .bib: \comumBibliografia (padrão: referencias.bib).
%             Para outro estilo, defina antes:
%               \newcommand{\comumEstiloBib}{authoryear}
% hyperref  : links no PDF (pretos para não "poluir" a impressão).
%             Ajuste título/autor com \hypersetup{pdftitle=..., pdfauthor=...}.
% cleveref  : \cref / \Cref em português, com nome capitalizado.
%             Precisa vir DEPOIS do hyperref — por isso fica neste módulo.
%
% Compilação: latexmk detecta o biber automaticamente.
% =====================================================================

\providecommand{\comumBibliografia}{referencias.bib}
\providecommand{\comumEstiloBib}{abnt}

\usepackage{csquotes}
\usepackage[backend=biber, style=\comumEstiloBib, sorting=nyt]{biblatex}
\addbibresource{\comumBibliografia}

\usepackage{hyperref}
\hypersetup{
  colorlinks=true,
  linkcolor=black,
  citecolor=black,
  urlcolor=blue,
}

\usepackage[brazilian, capitalise]{cleveref}

% =====================================================================
% comum/theoremstyles.tex — ambientes de teorema (amsthm) + cleveref
% =====================================================================
% Requer: amsthm (pacotes/matematica.tex) e cleveref (pacotes/referencias.tex).
%
% Ambientes disponíveis:
%   estilo plain      : theorem, proposition, lemma, corollary
%   estilo definition : definition, axiom, example, examples, observation
%   sem número        : theorem*, definition*, example*, observation*
%
% Configuração (defina ANTES de carregar este arquivo):
%   \newcommand{\comumNumeracao}{chapter}  % Teorema 2.1  (livro/TCC)
%   \newcommand{\comumNumeracao}{section}  % Teorema 2.1  (artigo)
%   (não definir)                           % Teorema 1, 2, ...
%
%   \newcommand{\comumContadorUnico}{}     % todos os ambientes dividem um
%                                          % contador (Def. 2.1, Teo. 2.2, ...)
%                                          % padrão: um contador por ambiente
%
% Uso com cleveref: \cref{teo:x} -> "teorema 2.1"; \Cref{teo:x} -> "Teorema 2.1".
% Prefixos de rótulo sugeridos: teo:, prop:, lem:, cor:, def:, ax:, ex:, obs:
% =====================================================================

\providecommand{\comumNumeracao}{}
\ifdefined\comumContadorUnico
  \RequirePackage{aliascnt}
\fi

\makeatletter
% \comum@novoteorema{ambiente}{Nome}
\newcommand{\comum@novoteorema}[2]{%
  \ifdefined\comumContadorUnico
    % o primeiro ambiente (theorem) cria o contador; os demais o reutilizam
    \@ifundefined{c@theorem}%
      {\comum@teoremasimples{#1}{#2}}%
      {% aliascnt: contador "apelido" de theorem, para o cleveref saber
       % que o rótulo é uma definição/lema/... e não um teorema
       \newaliascnt{#1}{theorem}%
       \newtheorem{#1}[#1]{#2}%
       \aliascntresetthe{#1}}%
  \else
    \comum@teoremasimples{#1}{#2}%
  \fi
}
\newcommand{\comum@teoremasimples}[2]{%
  \if\relax\detokenize\expandafter{\comumNumeracao}\relax
    \newtheorem{#1}{#2}%
  \else
    \newtheorem{#1}{#2}[\comumNumeracao]%
  \fi
}

% ---------- Estilo plain (itálico): resultados ----------
\theoremstyle{plain}
\comum@novoteorema{theorem}{Teorema}
\comum@novoteorema{proposition}{Proposição}
\comum@novoteorema{lemma}{Lema}
\comum@novoteorema{corollary}{Corolário}
\newtheorem*{theorem*}{Teorema}

% ---------- Estilo definition (reto): definições e exemplos ----------
\theoremstyle{definition}
\comum@novoteorema{definition}{Definição}
\comum@novoteorema{axiom}{Axioma}
\comum@novoteorema{example}{Exemplo}
% examples compartilha o contador de example
\newtheorem{examples}[example]{Exemplos}
\comum@novoteorema{observation}{Observação}
\newtheorem*{definition*}{Definição}
\newtheorem*{example*}{Exemplo}
\newtheorem*{observation*}{Observação}

% ---------- Fim de demonstração ----------
\renewcommand{\qedsymbol}{$\blacksquare$}

% ---------- Nomes para \cref (minúsculo) e \Cref (maiúsculo) ----------
\crefname{theorem}{teorema}{teoremas}
\Crefname{theorem}{Teorema}{Teoremas}
\crefname{proposition}{proposição}{proposições}
\Crefname{proposition}{Proposição}{Proposições}
\crefname{lemma}{lema}{lemas}
\Crefname{lemma}{Lema}{Lemas}
\crefname{corollary}{corolário}{corolários}
\Crefname{corollary}{Corolário}{Corolários}
\crefname{definition}{definição}{definições}
\Crefname{definition}{Definição}{Definições}
\crefname{axiom}{axioma}{axiomas}
\Crefname{axiom}{Axioma}{Axiomas}
\crefname{example}{exemplo}{exemplos}
\Crefname{example}{Exemplo}{Exemplos}
\crefname{examples}{exemplos}{exemplos}
\Crefname{examples}{Exemplos}{Exemplos}
\crefname{observation}{observação}{observações}
\Crefname{observation}{Observação}{Observações}
\makeatother

% =====================================================================
% comum/comandos.tex — macros e ambientes próprios
% =====================================================================
% \asyimage[opções]{arquivo}{Legenda}[rótulo]
%     Figura flutuante [H] centralizada com legenda e rótulo opcional.
%     Feita para os PDFs gerados a partir de .asy (TCC), mas serve para
%     qualquer imagem. Ex.:
%       \asyimage[width=0.6\linewidth]{asyimages/diagram1.pdf}
%                {Parametrizações}[fig:param]
%
% \begin{citacao} ... \end{citacao}
%     Citação longa no padrão ABNT (NBR 10520): recuo de 4cm,
%     fonte menor e espaçamento simples.
% =====================================================================

\NewDocumentCommand{\asyimage}{O{width=0.8\linewidth} m m o}{%
  \begin{figure}[H]%
    \centering
    \includegraphics[#1]{#2}%
    \caption{#3}%
    \IfValueT{#4}{\label{#4}}%
  \end{figure}%
}

\newenvironment{citacao}{%
  \par\smallskip
  \begingroup
  \leftskip=4cm
  \small
  \ifdefined\singlespacing\singlespacing\fi
  \noindent\ignorespaces
}{%
  \par\endgroup
  \smallskip
}


%
% Carrega, NESTA ORDEM, os módulos abaixo. A ordem importa:
%   1. idioma       — codificação, fontes e babel (sempre primeiro)
%   2. matematica   — ams*, mathtools, bm, tikz-cd
%   3. figuras      — graphicx, tikz, pgfplots, float, (sub)caption
%   4. layout-abnt  — margens ABNT, espaçamento 1,5, cabeçalho
%   5. referencias  — biblatex-abnt, hyperref e cleveref
%                     (hyperref/cleveref precisam vir depois de quase tudo)
%   6. theoremstyles — ambientes de teorema + nomes do cleveref
%                     (precisa do cleveref já carregado)
%   7. comandos     — macros próprias (\asyimage, ambiente citacao)
%
% Para desligar um módulo, comente a linha correspondente.
% Em beamer NÃO use este arquivo inteiro: beamer já define teoremas e
% layout; carregue só os módulos 1, 2, 3 e 5 (veja templates/slides-unespar).
% =====================================================================

% =====================================================================
% comum/pacotes/idioma.tex — codificação, fontes e idioma
% =====================================================================
% Detecta o motor de compilação:
%   * XeLaTeX / LuaLaTeX (padrão do repositório, ver .vscode/settings.json):
%       fontspec com Latin Modern (Unicode nativo, sem inputenc).
%   * pdfLaTeX: inputenc utf8 + fontenc T1 + lmodern.
% Para trocar a fonte principal em XeLaTeX, use DEPOIS deste módulo:
%       \setmainfont{Liberation Serif}   % ex.: Times-compatível (IC/PRPPG)
%       \setsansfont{texgyreheros}[Extension=.otf, ...]  % Helvetica (slides)
% =====================================================================

\usepackage{iftex}

\ifPDFTeX
  \usepackage[utf8]{inputenc}
  \usepackage[T1]{fontenc}
  \usepackage{lmodern}
\else
  \usepackage{fontspec}
\fi

\usepackage[brazilian]{babel}

% Listas: enumitem (rótulos/espaçamento) e moreenum (\roman, \greek, ...)
\usepackage{enumitem}
\usepackage{moreenum}

% =====================================================================
% comum/pacotes/matematica.tex — pacotes de matemática
% =====================================================================
% amsmath/amssymb/amsthm : base AMS (amsthm é usado por theoremstyles.tex)
% mathtools              : extensões do amsmath (\coloneqq, \DeclarePairedDelimiter)
% bm                     : negrito em símbolos matemáticos (\bm{\alpha})
% amscd / tikz-cd        : diagramas comutativos (simples / completos)
% =====================================================================

\usepackage{amsmath}
\usepackage{amssymb}
\usepackage{amsthm}
\usepackage{mathtools}
\usepackage{bm}
\usepackage{amscd}
\usepackage{tikz-cd}

% =====================================================================
% comum/pacotes/figuras.tex — figuras, TikZ e legendas
% =====================================================================
% graphicx            : \includegraphics (imagens ficam em ../figures/)
% tikz + bibliotecas  : as usadas nos projetos (setas, ângulos, sombras...)
% pgfplots            : gráficos de funções
% float               : opção [H] para fixar figuras
% caption/subcaption  : legendas e subfiguras
% asymptote           : figuras 3D em .asy (TCC). Desligado por padrão
%                       porque exige o programa `asy` instalado; para
%                       ligar, defina \newcommand{\comumAsymptote}{} antes
%                       de carregar os módulos (veja comandos.tex: \asyimage).
% =====================================================================

\usepackage{graphicx}
\graphicspath{{../figures/}{figures/}}

\usepackage{tikz}
\usetikzlibrary{arrows.meta, calc, positioning, decorations.markings,
  angles, quotes, shapes.geometric, shapes.symbols, shadows}

\usepackage{pgfplots}
\pgfplotsset{compat=1.18}

\usepackage{float}
\usepackage{caption}
\usepackage{subcaption}

\ifdefined\comumAsymptote
  \usepackage{asymptote}
\fi

% =====================================================================
% comum/pacotes/layout-abnt.tex — página no padrão ABNT (NBR 14724)
% =====================================================================
% Margens: superior/esquerda 3cm, inferior/direita 2cm.
% Espaçamento 1,5; recuo no primeiro parágrafo; número da página no
% canto superior direito (fancyhdr); tabelas com booktabs/multirow.
%
% Documentos com cabeçalho/margens próprios (ex.: relatório PRPPG)
% podem sobrescrever depois com \geometry{...} e \fancyhead{...}.
% NÃO usar em beamer.
% =====================================================================

\usepackage[a4paper, top=3cm, bottom=2cm, left=3cm, right=2cm,
  headheight=14.5pt]{geometry}

\usepackage{setspace}
\onehalfspacing

\usepackage{indentfirst}
\setlength{\parindent}{1.25cm}

\usepackage{fancyhdr}
\pagestyle{fancy}
\fancyhf{}
\fancyhead[R]{\thepage}
\renewcommand{\headrulewidth}{0.4pt}

\usepackage{booktabs}
\usepackage{multirow}

% =====================================================================
% comum/pacotes/referencias.tex — bibliografia e referências cruzadas
% =====================================================================
% biblatex  : backend biber, estilo ABNT (pacote biblatex-abnt).
%             Arquivo .bib: \comumBibliografia (padrão: referencias.bib).
%             Para outro estilo, defina antes:
%               \newcommand{\comumEstiloBib}{authoryear}
% hyperref  : links no PDF (pretos para não "poluir" a impressão).
%             Ajuste título/autor com \hypersetup{pdftitle=..., pdfauthor=...}.
% cleveref  : \cref / \Cref em português, com nome capitalizado.
%             Precisa vir DEPOIS do hyperref — por isso fica neste módulo.
%
% Compilação: latexmk detecta o biber automaticamente.
% =====================================================================

\providecommand{\comumBibliografia}{referencias.bib}
\providecommand{\comumEstiloBib}{abnt}

\usepackage{csquotes}
\usepackage[backend=biber, style=\comumEstiloBib, sorting=nyt]{biblatex}
\addbibresource{\comumBibliografia}

\usepackage{hyperref}
\hypersetup{
  colorlinks=true,
  linkcolor=black,
  citecolor=black,
  urlcolor=blue,
}

\usepackage[brazilian, capitalise]{cleveref}

% =====================================================================
% comum/theoremstyles.tex — ambientes de teorema (amsthm) + cleveref
% =====================================================================
% Requer: amsthm (pacotes/matematica.tex) e cleveref (pacotes/referencias.tex).
%
% Ambientes disponíveis:
%   estilo plain      : theorem, proposition, lemma, corollary
%   estilo definition : definition, axiom, example, examples, observation
%   sem número        : theorem*, definition*, example*, observation*
%
% Configuração (defina ANTES de carregar este arquivo):
%   \newcommand{\comumNumeracao}{chapter}  % Teorema 2.1  (livro/TCC)
%   \newcommand{\comumNumeracao}{section}  % Teorema 2.1  (artigo)
%   (não definir)                           % Teorema 1, 2, ...
%
%   \newcommand{\comumContadorUnico}{}     % todos os ambientes dividem um
%                                          % contador (Def. 2.1, Teo. 2.2, ...)
%                                          % padrão: um contador por ambiente
%
% Uso com cleveref: \cref{teo:x} -> "teorema 2.1"; \Cref{teo:x} -> "Teorema 2.1".
% Prefixos de rótulo sugeridos: teo:, prop:, lem:, cor:, def:, ax:, ex:, obs:
% =====================================================================

\providecommand{\comumNumeracao}{}
\ifdefined\comumContadorUnico
  \RequirePackage{aliascnt}
\fi

\makeatletter
% \comum@novoteorema{ambiente}{Nome}
\newcommand{\comum@novoteorema}[2]{%
  \ifdefined\comumContadorUnico
    % o primeiro ambiente (theorem) cria o contador; os demais o reutilizam
    \@ifundefined{c@theorem}%
      {\comum@teoremasimples{#1}{#2}}%
      {% aliascnt: contador "apelido" de theorem, para o cleveref saber
       % que o rótulo é uma definição/lema/... e não um teorema
       \newaliascnt{#1}{theorem}%
       \newtheorem{#1}[#1]{#2}%
       \aliascntresetthe{#1}}%
  \else
    \comum@teoremasimples{#1}{#2}%
  \fi
}
\newcommand{\comum@teoremasimples}[2]{%
  \if\relax\detokenize\expandafter{\comumNumeracao}\relax
    \newtheorem{#1}{#2}%
  \else
    \newtheorem{#1}{#2}[\comumNumeracao]%
  \fi
}

% ---------- Estilo plain (itálico): resultados ----------
\theoremstyle{plain}
\comum@novoteorema{theorem}{Teorema}
\comum@novoteorema{proposition}{Proposição}
\comum@novoteorema{lemma}{Lema}
\comum@novoteorema{corollary}{Corolário}
\newtheorem*{theorem*}{Teorema}

% ---------- Estilo definition (reto): definições e exemplos ----------
\theoremstyle{definition}
\comum@novoteorema{definition}{Definição}
\comum@novoteorema{axiom}{Axioma}
\comum@novoteorema{example}{Exemplo}
% examples compartilha o contador de example
\newtheorem{examples}[example]{Exemplos}
\comum@novoteorema{observation}{Observação}
\newtheorem*{definition*}{Definição}
\newtheorem*{example*}{Exemplo}
\newtheorem*{observation*}{Observação}

% ---------- Fim de demonstração ----------
\renewcommand{\qedsymbol}{$\blacksquare$}

% ---------- Nomes para \cref (minúsculo) e \Cref (maiúsculo) ----------
\crefname{theorem}{teorema}{teoremas}
\Crefname{theorem}{Teorema}{Teoremas}
\crefname{proposition}{proposição}{proposições}
\Crefname{proposition}{Proposição}{Proposições}
\crefname{lemma}{lema}{lemas}
\Crefname{lemma}{Lema}{Lemas}
\crefname{corollary}{corolário}{corolários}
\Crefname{corollary}{Corolário}{Corolários}
\crefname{definition}{definição}{definições}
\Crefname{definition}{Definição}{Definições}
\crefname{axiom}{axioma}{axiomas}
\Crefname{axiom}{Axioma}{Axiomas}
\crefname{example}{exemplo}{exemplos}
\Crefname{example}{Exemplo}{Exemplos}
\crefname{examples}{exemplos}{exemplos}
\Crefname{examples}{Exemplos}{Exemplos}
\crefname{observation}{observação}{observações}
\Crefname{observation}{Observação}{Observações}
\makeatother

% =====================================================================
% comum/comandos.tex — macros e ambientes próprios
% =====================================================================
% \asyimage[opções]{arquivo}{Legenda}[rótulo]
%     Figura flutuante [H] centralizada com legenda e rótulo opcional.
%     Feita para os PDFs gerados a partir de .asy (TCC), mas serve para
%     qualquer imagem. Ex.:
%       \asyimage[width=0.6\linewidth]{asyimages/diagram1.pdf}
%                {Parametrizações}[fig:param]
%
% \begin{citacao} ... \end{citacao}
%     Citação longa no padrão ABNT (NBR 10520): recuo de 4cm,
%     fonte menor e espaçamento simples.
% =====================================================================

\NewDocumentCommand{\asyimage}{O{width=0.8\linewidth} m m o}{%
  \begin{figure}[H]%
    \centering
    \includegraphics[#1]{#2}%
    \caption{#3}%
    \IfValueT{#4}{\label{#4}}%
  \end{figure}%
}

\newenvironment{citacao}{%
  \par\smallskip
  \begingroup
  \leftskip=4cm
  \small
  \ifdefined\singlespacing\singlespacing\fi
  \noindent\ignorespaces
}{%
  \par\endgroup
  \smallskip
}


%
% Carrega, NESTA ORDEM, os módulos abaixo. A ordem importa:
%   1. idioma       — codificação, fontes e babel (sempre primeiro)
%   2. matematica   — ams*, mathtools, bm, tikz-cd
%   3. figuras      — graphicx, tikz, pgfplots, float, (sub)caption
%   4. layout-abnt  — margens ABNT, espaçamento 1,5, cabeçalho
%   5. referencias  — biblatex-abnt, hyperref e cleveref
%                     (hyperref/cleveref precisam vir depois de quase tudo)
%   6. theoremstyles — ambientes de teorema + nomes do cleveref
%                     (precisa do cleveref já carregado)
%   7. comandos     — macros próprias (\asyimage, ambiente citacao)
%
% Para desligar um módulo, comente a linha correspondente.
% Em beamer NÃO use este arquivo inteiro: beamer já define teoremas e
% layout; carregue só os módulos 1, 2, 3 e 5 (veja templates/slides-unespar).
% =====================================================================

% =====================================================================
% comum/pacotes/idioma.tex — codificação, fontes e idioma
% =====================================================================
% Detecta o motor de compilação:
%   * XeLaTeX / LuaLaTeX (padrão do repositório, ver .vscode/settings.json):
%       fontspec com Latin Modern (Unicode nativo, sem inputenc).
%   * pdfLaTeX: inputenc utf8 + fontenc T1 + lmodern.
% Para trocar a fonte principal em XeLaTeX, use DEPOIS deste módulo:
%       \setmainfont{Liberation Serif}   % ex.: Times-compatível (IC/PRPPG)
%       \setsansfont{texgyreheros}[Extension=.otf, ...]  % Helvetica (slides)
% =====================================================================

\usepackage{iftex}

\ifPDFTeX
  \usepackage[utf8]{inputenc}
  \usepackage[T1]{fontenc}
  \usepackage{lmodern}
\else
  \usepackage{fontspec}
\fi

\usepackage[brazilian]{babel}

% Listas: enumitem (rótulos/espaçamento) e moreenum (\roman, \greek, ...)
\usepackage{enumitem}
\usepackage{moreenum}

% =====================================================================
% comum/pacotes/matematica.tex — pacotes de matemática
% =====================================================================
% amsmath/amssymb/amsthm : base AMS (amsthm é usado por theoremstyles.tex)
% mathtools              : extensões do amsmath (\coloneqq, \DeclarePairedDelimiter)
% bm                     : negrito em símbolos matemáticos (\bm{\alpha})
% amscd / tikz-cd        : diagramas comutativos (simples / completos)
% =====================================================================

\usepackage{amsmath}
\usepackage{amssymb}
\usepackage{amsthm}
\usepackage{mathtools}
\usepackage{bm}
\usepackage{amscd}
\usepackage{tikz-cd}

% =====================================================================
% comum/pacotes/figuras.tex — figuras, TikZ e legendas
% =====================================================================
% graphicx            : \includegraphics (imagens ficam em ../figures/)
% tikz + bibliotecas  : as usadas nos projetos (setas, ângulos, sombras...)
% pgfplots            : gráficos de funções
% float               : opção [H] para fixar figuras
% caption/subcaption  : legendas e subfiguras
% asymptote           : figuras 3D em .asy (TCC). Desligado por padrão
%                       porque exige o programa `asy` instalado; para
%                       ligar, defina \newcommand{\comumAsymptote}{} antes
%                       de carregar os módulos (veja comandos.tex: \asyimage).
% =====================================================================

\usepackage{graphicx}
\graphicspath{{../figures/}{figures/}}

\usepackage{tikz}
\usetikzlibrary{arrows.meta, calc, positioning, decorations.markings,
  angles, quotes, shapes.geometric, shapes.symbols, shadows}

\usepackage{pgfplots}
\pgfplotsset{compat=1.18}

\usepackage{float}
\usepackage{caption}
\usepackage{subcaption}

\ifdefined\comumAsymptote
  \usepackage{asymptote}
\fi

% =====================================================================
% comum/pacotes/layout-abnt.tex — página no padrão ABNT (NBR 14724)
% =====================================================================
% Margens: superior/esquerda 3cm, inferior/direita 2cm.
% Espaçamento 1,5; recuo no primeiro parágrafo; número da página no
% canto superior direito (fancyhdr); tabelas com booktabs/multirow.
%
% Documentos com cabeçalho/margens próprios (ex.: relatório PRPPG)
% podem sobrescrever depois com \geometry{...} e \fancyhead{...}.
% NÃO usar em beamer.
% =====================================================================

\usepackage[a4paper, top=3cm, bottom=2cm, left=3cm, right=2cm,
  headheight=14.5pt]{geometry}

\usepackage{setspace}
\onehalfspacing

\usepackage{indentfirst}
\setlength{\parindent}{1.25cm}

\usepackage{fancyhdr}
\pagestyle{fancy}
\fancyhf{}
\fancyhead[R]{\thepage}
\renewcommand{\headrulewidth}{0.4pt}

\usepackage{booktabs}
\usepackage{multirow}

% =====================================================================
% comum/pacotes/referencias.tex — bibliografia e referências cruzadas
% =====================================================================
% biblatex  : backend biber, estilo ABNT (pacote biblatex-abnt).
%             Arquivo .bib: \comumBibliografia (padrão: referencias.bib).
%             Para outro estilo, defina antes:
%               \newcommand{\comumEstiloBib}{authoryear}
% hyperref  : links no PDF (pretos para não "poluir" a impressão).
%             Ajuste título/autor com \hypersetup{pdftitle=..., pdfauthor=...}.
% cleveref  : \cref / \Cref em português, com nome capitalizado.
%             Precisa vir DEPOIS do hyperref — por isso fica neste módulo.
%
% Compilação: latexmk detecta o biber automaticamente.
% =====================================================================

\providecommand{\comumBibliografia}{referencias.bib}
\providecommand{\comumEstiloBib}{abnt}

\usepackage{csquotes}
\usepackage[backend=biber, style=\comumEstiloBib, sorting=nyt]{biblatex}
\addbibresource{\comumBibliografia}

\usepackage{hyperref}
\hypersetup{
  colorlinks=true,
  linkcolor=black,
  citecolor=black,
  urlcolor=blue,
}

\usepackage[brazilian, capitalise]{cleveref}

% =====================================================================
% comum/theoremstyles.tex — ambientes de teorema (amsthm) + cleveref
% =====================================================================
% Requer: amsthm (pacotes/matematica.tex) e cleveref (pacotes/referencias.tex).
%
% Ambientes disponíveis:
%   estilo plain      : theorem, proposition, lemma, corollary
%   estilo definition : definition, axiom, example, examples, observation
%   sem número        : theorem*, definition*, example*, observation*
%
% Configuração (defina ANTES de carregar este arquivo):
%   \newcommand{\comumNumeracao}{chapter}  % Teorema 2.1  (livro/TCC)
%   \newcommand{\comumNumeracao}{section}  % Teorema 2.1  (artigo)
%   (não definir)                           % Teorema 1, 2, ...
%
%   \newcommand{\comumContadorUnico}{}     % todos os ambientes dividem um
%                                          % contador (Def. 2.1, Teo. 2.2, ...)
%                                          % padrão: um contador por ambiente
%
% Uso com cleveref: \cref{teo:x} -> "teorema 2.1"; \Cref{teo:x} -> "Teorema 2.1".
% Prefixos de rótulo sugeridos: teo:, prop:, lem:, cor:, def:, ax:, ex:, obs:
% =====================================================================

\providecommand{\comumNumeracao}{}
\ifdefined\comumContadorUnico
  \RequirePackage{aliascnt}
\fi

\makeatletter
% \comum@novoteorema{ambiente}{Nome}
\newcommand{\comum@novoteorema}[2]{%
  \ifdefined\comumContadorUnico
    % o primeiro ambiente (theorem) cria o contador; os demais o reutilizam
    \@ifundefined{c@theorem}%
      {\comum@teoremasimples{#1}{#2}}%
      {% aliascnt: contador "apelido" de theorem, para o cleveref saber
       % que o rótulo é uma definição/lema/... e não um teorema
       \newaliascnt{#1}{theorem}%
       \newtheorem{#1}[#1]{#2}%
       \aliascntresetthe{#1}}%
  \else
    \comum@teoremasimples{#1}{#2}%
  \fi
}
\newcommand{\comum@teoremasimples}[2]{%
  \if\relax\detokenize\expandafter{\comumNumeracao}\relax
    \newtheorem{#1}{#2}%
  \else
    \newtheorem{#1}{#2}[\comumNumeracao]%
  \fi
}

% ---------- Estilo plain (itálico): resultados ----------
\theoremstyle{plain}
\comum@novoteorema{theorem}{Teorema}
\comum@novoteorema{proposition}{Proposição}
\comum@novoteorema{lemma}{Lema}
\comum@novoteorema{corollary}{Corolário}
\newtheorem*{theorem*}{Teorema}

% ---------- Estilo definition (reto): definições e exemplos ----------
\theoremstyle{definition}
\comum@novoteorema{definition}{Definição}
\comum@novoteorema{axiom}{Axioma}
\comum@novoteorema{example}{Exemplo}
% examples compartilha o contador de example
\newtheorem{examples}[example]{Exemplos}
\comum@novoteorema{observation}{Observação}
\newtheorem*{definition*}{Definição}
\newtheorem*{example*}{Exemplo}
\newtheorem*{observation*}{Observação}

% ---------- Fim de demonstração ----------
\renewcommand{\qedsymbol}{$\blacksquare$}

% ---------- Nomes para \cref (minúsculo) e \Cref (maiúsculo) ----------
\crefname{theorem}{teorema}{teoremas}
\Crefname{theorem}{Teorema}{Teoremas}
\crefname{proposition}{proposição}{proposições}
\Crefname{proposition}{Proposição}{Proposições}
\crefname{lemma}{lema}{lemas}
\Crefname{lemma}{Lema}{Lemas}
\crefname{corollary}{corolário}{corolários}
\Crefname{corollary}{Corolário}{Corolários}
\crefname{definition}{definição}{definições}
\Crefname{definition}{Definição}{Definições}
\crefname{axiom}{axioma}{axiomas}
\Crefname{axiom}{Axioma}{Axiomas}
\crefname{example}{exemplo}{exemplos}
\Crefname{example}{Exemplo}{Exemplos}
\crefname{examples}{exemplos}{exemplos}
\Crefname{examples}{Exemplos}{Exemplos}
\crefname{observation}{observação}{observações}
\Crefname{observation}{Observação}{Observações}
\makeatother

% =====================================================================
% comum/comandos.tex — macros e ambientes próprios
% =====================================================================
% \asyimage[opções]{arquivo}{Legenda}[rótulo]
%     Figura flutuante [H] centralizada com legenda e rótulo opcional.
%     Feita para os PDFs gerados a partir de .asy (TCC), mas serve para
%     qualquer imagem. Ex.:
%       \asyimage[width=0.6\linewidth]{asyimages/diagram1.pdf}
%                {Parametrizações}[fig:param]
%
% \begin{citacao} ... \end{citacao}
%     Citação longa no padrão ABNT (NBR 10520): recuo de 4cm,
%     fonte menor e espaçamento simples.
% =====================================================================

\NewDocumentCommand{\asyimage}{O{width=0.8\linewidth} m m o}{%
  \begin{figure}[H]%
    \centering
    \includegraphics[#1]{#2}%
    \caption{#3}%
    \IfValueT{#4}{\label{#4}}%
  \end{figure}%
}

\newenvironment{citacao}{%
  \par\smallskip
  \begingroup
  \leftskip=4cm
  \small
  \ifdefined\singlespacing\singlespacing\fi
  \noindent\ignorespaces
}{%
  \par\endgroup
  \smallskip
}

